% ====================================================================
\section{The simulation theorem}\label{sec:simulation}
% ====================================================================

\subsection{Statement}

\begin{theorem}[Simulation theorem]\label{thm:simulation}
There is a logspace transducer which, given an MSLP $\mathcal P$ in normal form with $m$
registers, outputs an explicit tree $T$, an initial configuration $x(0)\in\{0,1\}^{V(T)}$, a vertex
$v^\ast$ and a nonnegative integer $\tau$ in binary, such that:
\begin{enumerate}[label=(\roman*)]
\item\label{sim:tree} $T$ is a tree, and every vertex of $T$ has degree $1$, $3$ or $4$;
\item\label{sim:size} $|V(T)|=O(m^{13})$;
\item\label{sim:time} $\tau\le 34m^2+79m$, and $\tau$ is given by the explicit formula of
  \cref{cor:tau};
\item\label{sim:correct} under synchronous majority, $x_{v^\ast}(\tau)=\val(\mathcal P)$.
\end{enumerate}
A second logspace transducer outputs $(T',x'(0),z)$ satisfying \ref{sim:tree} and \ref{sim:size}, with
$x'_z(0)=0$ and
\[
\exists\, t\ge 1:\ x'_z(t)=1\quad\Longleftrightarrow\quad\val(\mathcal P)=1.
\]
By \cref{lem:normalform} and closure of logspace reductions under composition~\cite{Jon75,LL76},
the same holds for every MSLP, with $m$ replaced by $r+L+1$.
\end{theorem}

The rest of this section explains how the proof is organised and proves the principle that joins
its two layers.

\subsection{Two layers: a controlled processor and autonomous clocks}

The tree $T$ consists of a \emph{processor} $P$ of $O(m)$ vertices, together with $O(m)$ private
\emph{generator} trees attached to it by single edges.
\begin{itemize}
\item The processor (\cref{sec:modules,sec:processor}) is analysed as a \emph{controlled tree}. Some
  of its vertices, the \emph{control slots}, are leaves whose bits are prescribed words. The
  remaining vertices follow the majority rule. All local operations and the complete program are
  proved correct in this forced setting.
\item Each control slot is then replaced by a \emph{robust generator} (\cref{sec:clocks}): a private
  tree whose root reproduces the prescribed word whatever the processor does. The substitution
  principle below shows that this replacement does not change the evolution of the processor.
\end{itemize}

\begin{definition}[Controlled tree]\label{def:controlled}
A \emph{controlled tree} is a finite tree $P$ together with a set $S\subseteq V(P)$ of leaves,
called \emph{control slots}, a horizon $\Hz\in\mathbb N$, and for every $\sigma\in S$ a word
$w_\sigma:[0,\Hz]\to\{0,1\}$. The unique neighbour of a slot $\sigma$ is its \emph{consumer}. Given
initial bits on $V(P)\setminus S$, the \emph{forced evolution} is the open evolution of
$V(P)\setminus S$ (\cref{def:open}) with boundary $S$ and boundary history $y_\sigma(t)=w_\sigma(t)$
for $t\le\Hz$, extended arbitrarily for $t>\Hz$.
\end{definition}

\begin{definition}[Robust generator]\label{def:robust}
Let $\Gamma$ be a finite rooted tree with root $\gamma$ and initial bits $x(0)$, and let
$h:\mathbb N\to\{0,1\}$. Attach to $\gamma$ one new external vertex $\epsilon$, the
\emph{consumer}. We say that $\Gamma$ is a \emph{robust generator of $h$} if, for \emph{every}
history $y:\mathbb N\to\{0,1\}$ of $\epsilon$, the open evolution of $V(\Gamma)$ with boundary
history $y$ satisfies $x_\gamma(t)=h(t)$ for all $t\ge 0$. Degrees in $\Gamma$ are counted including
the edge $\gamma\epsilon$. The consumer is not counted in $|\Gamma|$.
\end{definition}

Robustness quantifies over all histories $y$. These include histories that depend on the earlier
states of $\Gamma$ itself, since every such feedback history is in particular one fixed function
$y:\mathbb N\to\{0,1\}$.

\begin{proposition}[Substitution principle]\label{prop:substitution}
Let $(P,S,\Hz,(w_\sigma))$ be a controlled tree. For every $\sigma\in S$ let $\Gamma_\sigma$ be a
robust generator of some $h_\sigma$ with $h_\sigma(t)=w_\sigma(t)$ for $0\le t\le\Hz$. Assume that the
$\Gamma_\sigma$ are pairwise disjoint and disjoint from $V(P)\setminus S$. Let $T$ be obtained from
$P$ by identifying each slot $\sigma$ with the root of $\Gamma_\sigma$, so that the consumer of the
slot becomes the consumer of $\Gamma_\sigma$. Then:
\begin{enumerate}[label=(\alph*)]
\item\label{sub:tree} $T$ is a tree. The degree of a vertex of $V(P)\setminus S$ is the same in
  $T$ as in $P$, and every vertex of $\Gamma_\sigma$ has in $T$ the degree it has in $\Gamma_\sigma$,
  consumer edge included.
\item\label{sub:same} For every common initial configuration of $V(P)\setminus S$ and of the
  generators, the trajectory of $T$ agrees on $V(P)\setminus S$ with the forced evolution at all
  times $0\le t\le\Hz+1$.
\end{enumerate}
\end{proposition}

\begin{proof}
\ref{sub:tree} Each slot is a leaf of $P$, and each $\Gamma_\sigma$ is a tree joined to
$V(P)\setminus S$ by exactly the one edge from its root to the consumer. Gluing disjoint trees along
single edges produces a tree. Degrees are preserved by construction.

\ref{sub:same} Let $x(\cdot)$ be the trajectory of $T$. The only vertex outside $\Gamma_\sigma$ that
is adjacent to $\Gamma_\sigma$ is the consumer $\epsilon_\sigma$. By \cref{lem:locality}, the
restriction of $x$ to $\Gamma_\sigma$ is the open evolution of $\Gamma_\sigma$ under the boundary
history $y(t)=x_{\epsilon_\sigma}(t)$. This $y$ is one particular function
$\mathbb N\to\{0,1\}$, so robustness gives $x_{\gamma_\sigma}(t)=h_\sigma(t)=w_\sigma(t)$ for
$t\le\Hz$. The boundary of $V(P)\setminus S$ in $T$ consists exactly of the roots
$\gamma_\sigma$. Their histories therefore agree with the prescribed words on $[0,\Hz]$, and
\cref{lem:locality} shows that $x$ agrees with the forced evolution on $V(P)\setminus S$ up to time
$\Hz+1$.
\end{proof}

\begin{remark}\label{rem:afterH}
After time $\Hz$ the generators follow their own histories $h_\sigma$, which need not be constant. In
our construction they alternate. The prediction query is made at time $\tau\le\Hz$, so this tail is
irrelevant there. For eventual activation the tail \emph{is} relevant, which is why
\cref{sec:ever} adds a reader whose verdict is immune to it.
\end{remark}

\subsection{Plan of the proof of \texorpdfstring{\cref{thm:simulation}}{the simulation theorem}}

\begin{enumerate}[label=(\arabic*)]
\item \Cref{sec:modules}: memory cells, junctions and the local primitives \prim{SET}, \prim{COPY},
  \prim{MOVE}, \prim{ROUTE}, \prim{GATE} and \prim{READ}, each with a contract in the forced setting.
\item \Cref{sec:clocks}: which words a single fixed-phase generator can produce (phase words); the
  phase calculus that checks every control word of every primitive; and the robust generators
  themselves.
\item \Cref{sec:processor}: the processor, its normal configurations, one program instruction as a
  composition of primitives through explicit milestones, background normalisation, and the
  forced-control simulation theorem (\cref{thm:forced}).
\item \Cref{sec:analysis}: substitution of generators, size, time, degree, acyclicity and logspace
  constructibility. This completes the proof of \cref{thm:simulation}, except for the
  eventual-activation variant.
\item \Cref{sec:completeness}: membership in $\Ptime$, and the eventual-activation variant.
\end{enumerate}
